\documentclass[xcolor=table,aspectratio=169]{beamer}
\usepackage{beamerthemesplit}
\usepackage{wrapfig}
\usetheme{SPbGU}
\usepackage{pdfpages}
\usepackage{amsmath}
\usepackage{cmap}
\usepackage[T2A]{fontenc}
\usepackage[utf8]{inputenc}
\usepackage[english]{babel}
\usepackage{indentfirst}
\usepackage{amsmath}
\usepackage{tikz}
\usepackage{multirow}
\usepackage[noend]{algpseudocode}
\usepackage{algorithm}
\usepackage{algorithmicx}
\usepackage{fancyvrb}
\usetikzlibrary{shapes,arrows}
%usepackage{fancyvrb}
%\usepackage{minted}
%\usepackage{verbments}
\usepackage{fontawesome}


\beamertemplatenavigationsymbolsempty

\title[ТФЯ и лин. ал. для анализа графов]{Теория формальных языков и линейная алгебра для анализа графов}
\subtitle{Летняя Школа 2021}
\institute[СПбГУ]{
JetBrains Research, Programming Languages and Tools Lab  \\
Санкт-Петербургский Государственный Университет
}

\author[Семён Григорьев]{Семён Григорьев}

\date{19.12.2020}

\begin{document}
{
\begin{frame}[fragile]
  \begin{tabular}{p{2.0cm} p{9.5cm} p{1cm}}
   \begin{center}
      \includegraphics[height=1.5cm]{pictures/jetbrainsResearch.pdf}
    \end{center}
    &
    \begin{center}
      \includegraphics[height=1.5cm]{pictures/YC_logo.pdf}
    \end{center}
    &
    \begin{center}
      \includegraphics[height=1.5cm]{pictures/SPbGU_Logo.png}
    \end{center}
  \end{tabular}
  \titlepage
\end{frame}
}


\begin{frame}[fragile]

  \frametitle{Руководитель}
  Семён Григорьев
\begin{itemize}
      \item Доцент кафедры информатики СПбГУ
      \item Руководитель группы в лаборатории языковых инструментов JetBrains Research
      \item Контакты
      \begin{itemize}
        \item s.v.grigoriev@spbu.ru
        \item semyon.grigorev@jetbrains.com
        \item rsdpisuy@gmail.com (он же телега, он же гуглочат)        
      \end{itemize}

\end{itemize}

\end{frame}

\begin{frame}[fragile]

  \frametitle{Область интересов}
  Линейная алгебра и теория формальных языков для высокопроизводительного анализа графов
\begin{itemize}
      \item GraphBLAS
      \begin{itemize}
        \item \href{https://graphblas.github.io/}{Домашняя страница}
        \item \href{https://github.com/GraphBLAS/GraphBLAS-Pointers}{Большая коллекция полезных и интересных ссылок}
      \end{itemize}
      \pause
      \item Поиск путей с ограничениями в терминах формальгных языков
      \begin{itemize}
        \item Статический анализ кода
        \item Запросы к графовым базам данных
        \item Знакомство с областью можно начать с \href{https://research.jetbrains.org/groups/plt_lab/projects/cfpq/}{наших результатов} или с \href{https://github.com/JetBrains-Research/FormalLanguageConstrainedReachability-LectureNotes}{заготовки конспекта лекций}
      \end{itemize}
      \pause
      \item Графовые базы данных
      \begin{itemize}
        \item Neo4j
        \item RedisGraph
        \item TigerGraph
        \item \href{https://db-engines.com/en/ranking/graph+dbms}{Больше примеров и рейтинг графовых БД}        
      \end{itemize}
\end{itemize}

\end{frame}


\begin{frame}[fragile]
\frametitle{О проекте}
\begin{itemize}

      \item Задчи
      \begin{itemize}
        \item Получение теоретических оценок сложности        
        \item Разработка алгоритмов
        \item Реализация алгоритмов
        \item Постановка экспериментов        
      \end{itemize}
      \pause
      \item Технологии
      \begin{itemize}
        \item VHDL, FPGA
        \item C, OpenCL C, GPGPU
        \item Python
        \item Haskell, Scala, Coq, F\#
      \end{itemize}
\end{itemize}    
\end{frame}


\begin{frame}[fragile]

  \frametitle{Репозитории (примеры)}
  \begin{itemize}

      \item \href{https://github.com/JetBrains-Research/spbla}{spbla (разреженная булева линейная алгебра)}
      \item \href{https://github.com/JetBrains-Research/cfpq_data}{CFPQ\_Data (набор данных для экспериментального исследования)}
      \item \href{https://github.com/JetBrains-Research/cfpq_pyalgo}{CFPQ\_PyAlgo (коллекция алгоритмов для поиска путей с ограничениями в виде контекстно-свободных грамматик)}
      
  \end{itemize}
\end{frame}



\end{document}
